\documentclass[english,hebrew]{article}
\usepackage[HE8,T1]{fontenc}
\usepackage[utf8]{inputenc}
\setcounter{secnumdepth}{-2}
\setlength{\parindent}{0bp}
\usepackage{calc}
\usepackage{amsmath}
\usepackage{amsthm}
\usepackage{amssymb}

\makeatletter
\@ifundefined{date}{}{\date{}}
\usepackage{graphicx}
\usepackage{geometry}
\newcommand{\limi}{\underset{n\rightarrow\infty}{\lim}}
\newcommand{\limxz}{\underset{x\rightarrow x_{0}}{\lim}}
\newcommand{\limfxz}{\underset{x\rightarrow x_{0}}{\lim}f(x)}
\newcommand{\limfxm}{\underset{x\rightarrow x_{0}^{-}}{\lim}f(x)}
\newcommand{\limfxp}{\underset{x\rightarrow x_{0}^{+}}{\lim}f(x)}
\newcommand{\limxm}{\underset{x\rightarrow x_{0}^{-}}{\lim}}
\newcommand{\limxp}{\underset{x\rightarrow x_{0}^{+}}{\lim}}
\newcommand{\sqrm}{M_{m\times m}(\mathbb{F})}
\newcommand{\seqa}{(a_{n})_{n=1}^{\infty}}
\newcommand{\seqx}{(x_{n})_{n=1}^{\infty}}
\newcommand{\funcdr}{f:D\rightarrow\mathbb{R}}
\newcommand{\der}{\limxz\frac{f(x)-f(x_{0})}{x-x_{0}}}
\usepackage[all]{pdfprivacy}

\makeatother

\usepackage{babel}
\begin{document}
\newgeometry{top=1cm, bottom=2cm, left=1cm, right=1cm}
\title{אלגברה ליניארית {\beginL 2\endL} - הרצאה {\beginL 5\endL}}
\author{נעם לב}
\date{\L{$25.10.2026$}}

\maketitle
\noindent\fbox{\begin{minipage}[t]{1\columnwidth - 2\fboxsep - 2\fboxrule}%
\begin{center}
האוניברסיטה העברית \textbar{} אלגברה ליניארית ){\beginL 2\endL}( -
\L{$80135$} \textbar{} סמסטר{\small{} }{\beginL 2027\endL}{\small{} }א'
)תשפ\textquotedblleft ז א'( \textbar{} מרצה: פרופסור יבגני סטרחוב
\textbar{} מתרגל: מר לזר גוטרמן
\par\end{center}%
\end{minipage}}
\begin{center}
הסיכום הוא לא רשמי ולא עבר אישור של הסגל, על כן אין להשתמש בו )ובין
השאר במספור שלו ובשמות הטענות שבו( כמקור רשמי. אם ישנן טעויות, הן
שלי ולא של הסגל.
\par\end{center}

\medskip{}

\tableofcontents{}

\newpage{}

\part{מנהלה}

מבנה המבחן יפורסם לפני המבחן. 

חלוקת המבחן הסופי: \L{$50\%$} מהציון - משפטים, הגדרות, חלק תיאורטי.

\L{$50\%$} - תרגילים.

הבחנים הם רק אימון ומשוב; תחשבו עליהם כמו על בונוס.

הספר של אקסלר קרוב מאוד לקורס ואפשר לקרוא בו.

\part{הפולינום המינימלי )המשך(}

התחלנו לדבר על הפולינום המינימלי של ה\textquotedblright ל.

כרגע עוד לא ברור איך הסיפור הזה קשור לבעיות ספקטרליות, כלומר מה הקשר
בין הפולינום המינימלי לערכים העצמיים של ה\textquotedblright ל.

היום נראה קשר כזה.

נתחיל בלמה על פולינומים: אפשר לחלק פולינום בפולינום נתון עם שארית,
כאשר דרגת השארית נמוכה מדרגת הפולינום הנתון.

אפשר לחשוב על זה כמו על תרגיל בלינארית 1.

\section{למה - חילוק עם שארית של פולינומים}

נתונים שני פולינומים מעל \L{$\mathbb{F}$}: \L{%
\[
p(x)=a_{n}x^{n}+\dots+a_{1}x+a_{0}
\]
}%
\L{%
\[
q(x)=b_{m}x^{m}+\dots+b_{1}x+b_{0}
\]
}%

כך ש- \L{%
\[
\forall0\le i\le n:a_{i}\in\mathbb{F},\quad\forall0\le i\le m:b_{i}\in\mathbb{F},\quad\underset{\deg\left(p\right)\ge\deg\left(q\right)}{\underbrace{n\ge m}},\quad a_{n},b_{m}\ne0_{\mathbb{F}}
\]
}%
אז קיימים פולינומים \L{$S(x),r(x)$} מעל \L{$\mathbb{F}$} כך ש- \L{%
\[
p(x)=S(x)q(x)+r(x)
\]
}%

וגם \L{%
\[
\deg(r)<\deg(q)=m
\]
}%

למעשה יש כאן שתי טענות.

הראשונה - שניתן לכתוב את \L{$p$} בצורה כזו; השנייה, הלא-טריוויאלית
- שהשארית \L{$r$} היא לא סתם פולינום אלא פולינום מדרגה נמוכה מדרגת
\L{$q$}.

אפשר היה להשאיר את זה בלי הוכחה, אבל יש הוכחה יפה בעזרת אלגברה לינארית.
\begin{proof}
1. נגדיר \L{%
\[
V=\mathcal{P}_{n}(\mathbb{F})=\left\{ p\::\:p\text{ is a polynomial over }\mathbb{F},\;0\leq\deg(p)\leq n\right\} 
\]
}%
זהו מ\textquotedblright ו מעל \L{$\mathbb{F}$}, ו- \L{%
\[
\dim\mathcal{P}_{n}(\mathbb{F})=n+1
\]
}%

כי הבסיס הסטנדרטי שלו הוא \L{$\left(1,x,\dots,x^{n}\right)$} בו \L{$n+1$}
איברים.

זה המרחב \L{$\mathbb{F}[x]_{\le n}$} מלינארית 1, הרצאה 12.

פולינום האפס שייך אליו )אחרת זה לא מ\textquotedblright ו( ונהוג להגדיר
\L{$\deg(0)=-\infty$}, כי זה משמר את כללי המעלה.

למשל, לכל פולינום \L{$a$} מתקיים \L{%
\[
\deg(a\cdot0)=\deg(0)=-\infty=\deg(a)+(-\infty)=\deg(a)+\deg(0)
\]
}%
\L{%
\[
\deg(a+0)=\deg(a)=\max\left\{ \deg(a),-\infty\right\} =\max\left\{ \deg(a),\deg(0)\right\} 
\]
}%
כלומר הכללים \L{%
\[
\deg(ab)=\deg(a)+\deg(b)
\]
}%
\L{%
\[
\deg(a+b)\le\max\left\{ \deg(a),\deg(b)\right\} 
\]
}%
נשארים נכונים גם כשאחד הפולינומים הוא פולינום האפס.

בנוסף, \L{$-\infty<m$} לכל \L{$m$}, ולכן התנאי \L{$\deg(r)<\deg(q)$}
בלמה מתקיים גם כשהשארית היא \L{$0$}.

2. נגדיר את הרשימה הסדורה הבאה: \L{%
\[
\mathcal{B}_{V}=\left(\underbrace{1,x,\dots,x^{m-1}}_{m\text{ elements}},\underbrace{q(x),xq(x),\dots,x^{n-m}q(x)}_{n-m+1\text{ elements}}\right)
\]
}%
\L{%
\[
m+(n-m+1)=n+1=\dim\mathcal{P}_{n}(\mathbb{F})
\]
}%
דרגות הוקטורים ברשימה הן \L{%
\[
0,1,\dots,m-1,m,m+1,\dots,n
\]
}%
 כי \L{%
\[
\deg\left(x^{k}q(x)\right)=m+k,\quad b_{m}\ne0_{\mathbb{F}}
\]
}%

כלומר, דרגת הוקטורים עולה ממש לאורך הרשימה.

צ\textquotedblright ל של פולינומים מדרגות קטנות ממש מ-\L{$d$} הוא
פולינום מדרגה קטנה ממש מ-\L{$d$}, לכן לא קיים וקטור ברשימה השייך
לפרוש של קודמיו.

לכן לפי למת התלות הלינארית מלינארית {\beginL 1\endL} הרצאה {\beginL 9\endL}
האומרת כי סדרה ת\textquotedblright ל אם\textquotedblright ם קיים בה
וקטור השייך לפרוש של הקודמים לו )שזו ה-\L{$\text{Linear Dependence Lemma}$},
{\beginL 2.19\endL} במהדורה הרביעית של אקסלר( הרשימה הזאת בת\textquotedblright ל.

אז זאת רשימה בת\textquotedblright ל באורך המימד של המ\textquotedblright ו,
לכן היא בסיס )משפט מלינארית {\beginL 1\endL} הרצאה {\beginL 21\endL}(. 

דהיינו, \L{$\mathcal{B}_{V}$} בסיס של \L{$V$}.

3. \L{$p(x)\in\mathcal{P}_{n}(\mathbb{F})$}, ולכן קיימים \L{$\alpha_{0},\dots,\alpha_{m-1},\beta_{0},\dots,\beta_{n-m}\in\mathbb{F}$}
כך ש- \L{%
\[
p(x)=\alpha_{0}+\alpha_{1}x+\dots+\alpha_{m-1}x^{m-1}+\beta_{0}q(x)+\beta_{1}xq(x)+\dots+\beta_{n-m}x^{n-m}q(x)
\]
}%
\L{%
\[
=\underbrace{\left(\beta_{0}+\beta_{1}x+\dots+\beta_{n-m}x^{n-m}\right)}_{S(x)}q(x)+\underbrace{\left(\alpha_{0}+\alpha_{1}x+\dots+\alpha_{m-1}x^{m-1}\right)}_{r(x)}
\]
}%
ולכן \L{%
\[
p(x)=S(x)q(x)+r(x)
\]
}%
\L{%
\[
\deg\left(r(x)\right)\leq m-1<m=\deg\left(q(x)\right)
\]
}%
\L{%
\[
\deg\left(S(x)\right)\le n-m\le n=\deg\left(p(x)\right)
\]
}%
\end{proof}

\section{משפט - הפולינום המינימלי מחלק כל פולינום שמאפס את \L{$T$}}

יהי \L{$V$} מ\textquotedblright ו מעל \L{$\mathbb{F}$} כך ש- \L{%
\[
\dim V=n\geq1,\quad T\in\mathcal{L}(V)
\]
}%
נסמן ב-\L{$q_{\min}$} את הפולינום המינימלי של \L{$T$}.

יהי \L{$p$} פולינום המאפס את \L{$T$}, כלומר \L{$p(T)=0$}.

אז \L{$q_{\min}$} מחלק את \L{$p$} )ללא שארית(, כלומר קיים פולינום
\L{$S$} כך ש- \L{%
\[
p(x)=S(x)q_{\min}(x)
\]
}%

\begin{proof}
1. מהגדרת הפולינום המינימלי, \L{$\deg\left(p\right)\ge\deg\left(q_{\text{min}}\right)$}.

לכן, לפי הלמה הקודמת קיימים פולינומים \L{$S,r$} כך ש- \L{%
\[
p(x)=S(x)q_{\min}(x)+r(x)
\]
}%
\L{%
\[
\deg(r)<\deg(q_{\min})
\]
}%
צ\textquotedblright ל: \L{$r=0$}.

2. \L{%
\[
r(x)=p(x)-S(x)q_{\min}(x)
\]
}%
נציב \L{$T$}: \L{%
\[
r(T)=p(T)-S(T)q_{\min}(T)=0-S(T)\cdot0=0
\]
}%

3. קיבלנו ש-\L{$r(T)=0$} וגם \L{$\deg(r)<\deg(q_{\min})$}. 

אבל \L{$q_{\min}$} הוא הפולינום המינימלי של \L{$T$}, כלומר אין פולינום
שונה מאפס מדרגה קטנה מ-\L{$\deg(q_{\min})$} שמאפס את \L{$T$}. לכן
\L{%
\[
r=0
\]
}%
\L{%
\[
\therefore\ p(x)=S(x)q_{\min}(x)
\]
}%
\end{proof}

\subsection{הערה - מדוע מותר להציב \L{$T$} במקום \L{$x$}}

בהוכחה עברנו משוויון של פולינומים לשוויון של אופרטורים: \L{%
\[
r(x)=p(x)-S(x)q_{\min}(x)\quad\Longrightarrow\quad r(T)=p(T)-S(T)q_{\min}(T)
\]
}%

מה מצדיק את המעבר: שוויון של שני פולינומים \L{$a=b$} פירושו שלכל
\L{$k$} המקדם של \L{$x^{k}$} ב-\L{$a$} שווה למקדם של \L{$x^{k}$}
ב-\L{$b$}.

אבל \L{$a(T)$} ו-\L{$b(T)$} הם אותם מקדמים כפול אותן חזקות \L{$T^{k}$},
ולכן \L{$a(T)=b(T)$}.

כדי לעבור גם מסכום ומכפלה של פולינומים לסכום ומכפלה של אופרטורים צריך
שלכל שני פולינומים \L{$a,b$} מעל \L{$\mathbb{F}$} יתקיים \L{%
\[
(a+b)(T)=a(T)+b(T),\qquad(ab)(T)=a(T)\,b(T)
\]
}%

השוויון הראשון נובע מכך שהמקדמים של \L{$a+b$} הם סכומי המקדמים. 

לגבי השוויון השני, אם פותחים סוגריים )כל איבר של הסכום הראשון כפול
כל איבר של הסכום השני( מקבלים, \L{%
\[
a(T)\,b(T)=\sum_{i=0}^{k}\sum_{j=0}^{l}\alpha_{i}\beta_{j}T^{i}T^{j}\underset{T^{i}T^{j}=T^{i+j}}{=}\sum_{i=0}^{k}\sum_{j=0}^{l}\alpha_{i}\beta_{j}T^{i+j}=(ab)(T)
\]
}%

ובאופן מפורש\L{%
\[
a(T)\,b(T)=\left(\alpha_{0}I+\alpha_{1}T+\dots+\alpha_{k}T^{k}\right)\left(\beta_{0}I+\beta_{1}T+\dots+\beta_{l}T^{l}\right)
\]
}%
\L{%
\begin{align*}
 & \underset{T^{i}T^{j}=T^{i+j}}{=}\alpha_{0}\beta_{0}I+\alpha_{0}\beta_{1}T+\alpha_{0}\beta_{2}T^{2}+\dots+\alpha_{0}\beta_{l}T^{l}\\
 & \quad+\alpha_{1}\beta_{0}T+\alpha_{1}\beta_{1}T^{2}+\alpha_{1}\beta_{2}T^{3}+\dots+\alpha_{1}\beta_{l}T^{l+1}\\
 & \quad\ \vdots\\
 & \quad+\alpha_{k}\beta_{0}T^{k}+\alpha_{k}\beta_{1}T^{k+1}+\alpha_{k}\beta_{2}T^{k+2}+\dots+\alpha_{k}\beta_{l}T^{k+l}
\end{align*}
}%
וכשאוספים לפי חזקות של \L{$T$}: \L{%
\[
=\alpha_{0}\beta_{0}I+\left(\alpha_{0}\beta_{1}+\alpha_{1}\beta_{0}\right)T+\left(\alpha_{0}\beta_{2}+\alpha_{1}\beta_{1}+\alpha_{2}\beta_{0}\right)T^{2}+\dots+\left(\alpha_{k-1}\beta_{l}+\alpha_{k}\beta_{l-1}\right)T^{k+l-1}+\alpha_{k}\beta_{l}T^{k+l}
\]
}%
המקדמים בסוגריים הם בדיוק המקדמים של הפולינום \L{$ab$}, כלומר זה
\L{$(ab)(T)$}.

איפה הסכנה: בפולינום במשתנה אחד הכפל חילופי )\L{$x^{i}x^{j}=x^{j}x^{i}$}(,
והצבה שומרת על זה כי כל החזקות של אותו אופרטור \L{$T$} מתחלפות זו
עם זו: \L{$T^{i}T^{j}=T^{i+j}=T^{j}T^{i}$}.

לכן מותר להציב אופרטור יחיד \L{$T$} )או מטריצה יחידה \L{$A$}( בכל
זהות פולינומית במשתנה אחד.

למשל מ- \L{$(x-c)(x+d)=x^{2}+(d-c)x-cd$} נובע \L{%
\[
(T-cI)(T+dI)=T^{2}+(d-c)T-cdI
\]
}%

לעומת זאת, בזהות עם שני משתנים אסור להציב שני אופרטורים \L{$A,B$}
שאינם מתחלפים: הזהות \L{$(x-c)(y+d)=(y+d)(x-c)$} נכונה לפולינומים,
אבל \L{%
\[
(A-cI)(B+dI)=AB+dA-cB-cdI\ne BA+dA-cB-cdI=(B+dI)(A-cI)
\]
}%
כאשר \L{$AB\ne BA$}.

\section{משפט - שורשי הפולינום המינימלי הם הערכים העצמיים}

זה הקשר בין \L{$q_{\text{min}}$} של \L{$T$} לבין הערכים העצמיים
של \L{$T$}.

יהי \L{$V$} מ\textquotedblright ו מעל \L{$\mathbb{F}$} כך ש- \L{%
\[
\dim V=n\geq1,\quad T\in\mathcal{L}(V)
\]
}%

יהי \L{$q_{\min}$} הפולינום המינימלי של \L{$T$}, ויהי \L{$\lambda\in\mathbb{F}$}.

אז שני התנאים הבאים שקולים:

)I( \L{%
\[
q_{\min}(\lambda)=0_{\mathbb{F}}
\]
}%

)II( \L{$\lambda$} - ערך עצמי של \L{$T$}.\\

המשפט נכון מעל כל שדה, לא רק \L{$\mathbb{R}$} או \L{$\mathbb{C}$},
ולכן הוא מעניין: שורשי הפולינום המינימלי נותנים מיד את כל הערכים העצמיים,
בדיוק כמו שורשי הפולינום האופייני.

זה המשפט הבסיסי הראשון על פולינומים מינימליים שצריך לדעת. 

הפולינומים עצמם לא מעניינים אותנו; מעניין איך משתמשים בהם כדי לומר
משהו על ה\textquotedblright ל, על ערכיה העצמיים ו-וקטוריה העצמיים,
ואיך מוצאים בסיס שבו המטריצה המייצגת פשוטה ככל האפשר.

טענות על פולינומים נוכיח כלמות, אבל הן לא המטרה.
\begin{proof}
)II( \L{$\Rightarrow$} )I(:

1. נכתוב \L{%
\[
q_{\min}(x)=x^{m}+b_{m-1}x^{m-1}+\dots+b_{1}x+b_{0}
\]
}%
לפי )I(, \L{$q_{\min}(\lambda)=0_{\mathbb{F}}$}, ולכן \L{%
\begin{align*}
q_{\min}(x) & =q_{\min}(x)-\underbrace{q_{\min}(\lambda)}_{=0_{\mathbb{F}}}\\
 & =\left(x^{m}-\lambda^{m}\right)+b_{m-1}\left(x^{m-1}-\lambda^{m-1}\right)+\dots+b_{1}\left(x-\lambda\right)
\end{align*}
}%
\L{%
\begin{align*}
 & \underset{(1)}{=}\left(x-\lambda\right)\left(x^{m-1}+x^{m-2}\lambda+\dots+x\lambda^{m-2}+\lambda^{m-1}\right)\\
 & \quad+b_{m-1}\left(x-\lambda\right)\left(x^{m-2}+x^{m-3}\lambda+\dots+x\lambda^{m-3}+\lambda^{m-2}\right)\\
 & \quad\quad\vdots\\
 & \quad+b_{2}\left(x-\lambda\right)\left(x+\lambda\right)\\
 & \quad+b_{1}\left(x-\lambda\right)\\
 & =\left(x-\lambda\right)\underset{S(x)}{\underbrace{\left[\left(x^{m-1}+\dots+\lambda^{m-1}\right)+b_{m-1}\left(x^{m-2}+\dots+\lambda^{m-2}\right)+\dots+b_{2}\left(x+\lambda\right)+b_{1}\right]}}
\end{align*}
}%
\L{$(1)$} לכל \L{$k\geq2$},\L{%
\[
x^{k}-\lambda^{k}=(x-\lambda)\left(x^{k-1}+x^{k-2}\lambda+\dots+x\lambda^{k-2}+\lambda^{k-1}\right)
\]
}%
למשל, \L{$x^{2}-\lambda^{2}=(x-\lambda)(x+\lambda)$}. 

לכן ניתן להוציא גורם משותף \L{$(x-\lambda)$} מכל המחוברים וכמו בפיתוח
נקבל \L{%
\[
q_{\min}(x)=(x-\lambda)S(x)
\]
}%
כאשר \L{$S(x)$} פולינום המקיים \L{%
\[
\deg S(x)=m-1<\deg q_{\min}
\]
}%

הפולינום המפורש \L{$S$} לא חשוב, רק העובדה שדרגתו נמוכה מדרגת הפולינום
המינימלי )המחובר הראשון תורם \L{$x^{m-1}$}, האחרים חזקות נמוכות יותר(.

2. נציב \L{$T$} )כיסינו בהערה מדוע זה מותר(: \L{%
\[
0\underset{q_{\text{min}}\text{ def}}{=}q_{\min}(T)=(T-\lambda I)S(T)
\]
}%
מצד שני, \L{%
\[
S(T)\neq0
\]
}%
כי \L{$S$} פולינום שונה מאפס עם \L{$\deg S<\deg q_{\min}$}, ואילו
\L{$S(T)=0$} היה \L{$S$} פולינום מדרגה נמוכה יותר שמאפס את \L{$T$},
בסתירה למינימליות של \L{$q_{\min}$}.

3. לכן קיים \L{$0_{V}\ne v_{0}\in V$} כך ש- \L{%
\[
v'\overset{\text{def}}{=}S(T)\left(v_{0}\right)\neq0_{V}
\]
}%
נפעיל את השוויון מ-2 על \L{$v_{0}$}\L{%
\[
\left(T-\lambda I\right)\left(S\left(T\right)\left(v_{0}\right)\right)=0\left(v_{0}\right)=0_{V}
\]
}%
כלומר\L{%
\[
T\left(S\left(T\right)\left(v_{0}\right)\right)-\lambda I\left(S\left(T\right)\left(v_{0}\right)\right)=0_{V}
\]
}%
\L{%
\[
\therefore\ T\left(S\left(T\right)\left(v_{0}\right)\right)=\lambda I\left(S\left(T\right)\left(v_{0}\right)\right)
\]
}%
\L{%
\[
\therefore T\left(v'\right)=\lambda v'
\]
}%

מכאן \L{$\lambda$} הוא ערך עצמי של \L{$T$}, ו-\L{$v'$} וקטור עצמי
השייך לו.

יבגני אמר: זה משפט נחמד לתת במבחן - הוא קל ויפה, אין בו יותר מדי.
יהיו משפטים יותר קשים.

)I( \L{$\Rightarrow$} )II(:

1. \L{$\lambda$} ערך עצמי, לכן קיים \L{$0_{V}\ne v\in V$} כך ש-
\L{$Tv=\lambda v$}.

2. נפעיל את \L{$T$} על שני הצדדים: \L{%
\begin{align*}
 & T^{2}v=T\left(T\left(v\right)\right)=T(\lambda v)=\lambda Tv=\lambda^{2}v\\
 & T^{3}v=\lambda^{3}v\\
 & \quad\quad\ \vdots\\
 & T^{m}v=\lambda^{m}v
\end{align*}
}%
3. נזכיר כי \L{%
\[
q_{\min}(x)=x^{m}+b_{m-1}x^{m-1}+\dots+b_{1}x+b_{0}
\]
}%
ולכן \L{%
\[
0_{\mathcal{L}(V)}=q_{\min}(T)=T^{m}+b_{m-1}T^{m-1}+\dots+b_{1}T+b_{0}I
\]
}%
נפעיל על \L{$v$}: \L{%
\begin{align*}
0_{V} & =q_{\min}(T)\left(v\right)=T^{m}\left(v\right)+b_{m-1}T^{m-1}\left(v\right)+\dots+b_{1}T\left(v\right)+b_{0}\left(v\right)\\
 & \overset{(2)}{=}\lambda^{m}v+b_{m-1}\lambda^{m-1}v+\dots+b_{1}\lambda v+b_{0}v\\
 & =\left(\lambda^{m}+b_{m-1}\lambda^{m-1}+\dots+b_{1}\lambda+b_{0}\right)v\\
 & =q_{\min}(\lambda)\cdot v
\end{align*}
}%
כלומר \L{%
\[
0_{V}=q_{\min}(\lambda)\cdot v
\]
}%
כאשר \L{$q_{\min}(\lambda)\in\mathbb{F}$} כלומר הוא סקלר ו-\L{$v\neq0_{V}$}.

כפל של וקטור שונה מאפס בסקלר נותן \L{$0_{V}$} רק אם הסקלר הוא \L{$0_{\mathbb{F}}$},
לכן \L{%
\[
q_{\min}(\lambda)=0_{\mathbb{F}}
\]
}%
\end{proof}
שימו לב להבדל בין האובייקטים:

\L{$q_{\min}(x)$} הוא פולינום

\L{$q_{\min}(T)$} הוא אופרטור )סכום של חזקות של \L{$T$} כפול סקלרים
- זו ה\textquotedblright ל על \L{$V$} )אופרטור(, וכאן זה אופרטור
האפס(

\L{$q_{\min}(T)\left(v\right)$} הוא וקטור

\L{$q_{\min}(\lambda)$} הוא סקלר\\

המשפט נותן קשר בין שני הדברים המהותיים שראינו - ערכים עצמיים והפולינום
המינימלי.

המוטיבציה: השאלות שנותרו פתוחות הן מתי קיים בסיס של וקטורים עצמיים
מעל שדה כללי, ומתי מגיעים למטריצה משולשית עליונה; מתברר שעל שאלות
אלה אפשר לענות בעזרת הפולינום המינימלי.

\section{משפט - המקדם החופשי של הפולינום המינימלי והפיכות \L{$T$}}

יהי \L{$V$} מ\textquotedblright ו מעל \L{$\mathbb{F}$} כך ש- \L{%
\[
\dim V=n\geq1,\quad T\in\mathcal{L}(V)
\]
}%
\L{%
\[
q_{\min}(x)=x^{m}+b_{m-1}x^{m-1}+\dots+b_{1}x+b_{0}
\]
}%

אז התנאים הבאים שקולים:

)I( \L{%
\[
b_{0}=0_{\mathbb{F}}
\]
}%

)II( \L{$T$} אינה חח\textquotedblright ע\\

מכאן גם אינה על ובפרט אינה הפיכה, כי עבור אופרטור על מ\textquotedblright ו
נ\textquotedblright ס זה שקול כפי שראינו בהרצאה {\beginL \L{$16$}\endL}
של לינארית {\beginL 1\endL}.

גרסה אחרת של המשפט היא ש-\L{$T$} הפיכה אם\textquotedblright ם \L{$b_{0}\neq0_{\mathbb{F}}$}.

כלומר, ברגע שיש פולינום מינימלי, מיד אפשר לומר משהו על ההעתקה עצמה
- האם היא הפיכה - וזו הייתה שאלה מרכזית בלינארית 1.
\begin{proof}
)II( \L{$\Rightarrow$} )I(:

1. \L{$b_{0}=0_{\mathbb{F}}$}, לכן \L{%
\[
q_{\min}(x)=x^{m}+b_{m-1}x^{m-1}+\dots+b_{1}x
\]
}%
\L{%
\[
\therefore\ q_{\min}(0_{\mathbb{F}})=0_{\mathbb{F}}
\]
}%
כלומר \L{$0_{\mathbb{F}}$} הוא שורש של \L{$q_{\min}$}.

2. לפי המשפט הקודם, \L{$\lambda=0_{\mathbb{F}}$} הוא ערך עצמי של
\L{$T$}.

אז קיים \L{$0_{V}\ne v_{0}\in V$} כך ש- \L{%
\[
T(v_{0})=0_{\mathbb{F}}\cdot v_{0}=0_{V}
\]
}%
\L{%
\[
\therefore v_{0}\in\text{Null}(T),\;v_{0}\neq0_{V}
\]
}%
אז \L{$\text{Null}(T)\neq\left\{ 0_{V}\right\} $}, ולכן \L{$T$}
אינה חח\textquotedblright ע )לינארית 1, הרצאה \L{$16$}(.

)I( \L{$\Rightarrow$} )II(:

1. \L{$T$} אינה חח\textquotedblright ע, לכן קיים \L{$0_{V}\ne v_{0}\in\text{Null}(T)$}.
אז מתקיים \L{%
\[
T(v_{0})=0_{V}=0_{\mathbb{F}}\cdot v_{0}
\]
}%
\L{$\Leftarrow$} \L{$\lambda=0_{\mathbb{F}}$} ערך עצמי של \L{$T$}.

2. לפי המשפט הקודם, \L{$0_{\mathbb{F}}$} הוא שורש של \L{$q_{\min}(x)$}:
\L{%
\[
0_{\mathbb{F}}=q_{\min}(0_{\mathbb{F}})=0^{m}+b_{m-1}0^{m-1}+\dots+b_{1}\cdot0+b_{0}=b_{0}
\]
}%
\L{%
\[
\therefore\ b_{0}=0_{\mathbb{F}}
\]
}%
\end{proof}
השלב הבא הוא להוכיח שאם מציבים את \L{$T$} בפולינום האופייני מקבלים
\L{$0$} - זה משפט קיילי-המילטון.

הוא נראה פשוט, אבל ההוכחה די קשה; המשפט לא יידרש להוכחה במבחן, אבל
יוכח בהרצאה. 

לשם כך צריך קודם למה טכנית על מטריצות - חישוב המטריצה ההופכית בעזרת
דטרמיננטות - שהיא חומר של לינארית 1.

ההמשך בהרצאה הבאה.

\part{לקראת משפט קיילי-המילטון}

\section{משפט - נוסחה למטריצה ההופכית באמצעות דטרמיננטות}

תהי \L{%
\[
A=\begin{pmatrix}a_{1,1} & a_{1,2} & \cdots & a_{1,n}\\
a_{2,1} & a_{2,2} & \cdots & a_{2,n}\\
\vdots & \vdots & \ddots & \vdots\\
a_{n,1} & a_{n,2} & \cdots & a_{n,n}
\end{pmatrix}\in\text{Mat}\left(n\times n,\mathbb{F}\right)\text{ s.t. }\det(A)\ne0_{\mathbb{F}}
\]
}%

נגדיר \L{%
\[
B=\begin{pmatrix}b_{1,1} & b_{1,2} & \cdots & b_{1,n}\\
b_{2,1} & b_{2,2} & \cdots & b_{2,n}\\
\vdots & \vdots & \ddots & \vdots\\
b_{n,1} & b_{n,2} & \cdots & b_{n,n}
\end{pmatrix}\in\text{Mat}\left(n\times n,\mathbb{F}\right)
\]
}%

כאשר לכל \L{$1\le i,j\le n$}

\L{%
\[
b_{i,j}=\frac{(-1)^{i+j}}{\det(A)}\det\left(\begin{array}{cc|c|cc}
a_{1,1} & \cdots & a_{1,i} & \cdots & a_{1,n}\\
\vdots &  & \vdots &  & \vdots\\
\hline a_{j,1} & \cdots & a_{j,i} & \cdots & a_{j,n}\\
\hline \vdots &  & \vdots &  & \vdots\\
a_{n,1} & \cdots & a_{n,i} & \cdots & a_{n,n}
\end{array}\right)
\]
}%

הקווים מסמנים את השורה ה-\L{$j$} ואת העמודה ה-\L{$i$}, שאותן מוחקים;
הדטרמיננטה היא של המטריצה \L{$(n-1)\times(n-1)$} שנותרת.

כלומר, \L{$b_{i,j}$} מתקבל ממחיקת השורה ה-\L{$j$} והעמודה ה-\L{$i$}
של \L{$A$} )שימו לב להיפוך האינדקסים(, חישוב הדטרמיננטה של המטריצה
שנותרה, כפל בסימן \L{$\left(-1\right)^{i+j}$} וחלוקה ב-\L{$\det(A)$}.

המטריצה שנותרה היא המינור \L{$A_{j,i}$} )לינארית 1, הרצאה \L{$28$}(,
ולכן \L{%
\[
b_{i,j}=\frac{\left(-1\right)^{i+j}\det\left(A_{j,i}\right)}{\det(A)}
\]
}%

אז\L{%
\[
AB=I_{n}
\]
}%
 

יבגני: אם אתם לא מכירים את הנוסחה, בואו נבדוק שהיא נכונה. 

מומלץ מאוד לבצע את החישוב בעצמכם - לא בגלל המבחן )זה חומר של לינארית
1 ולא יישאל(, אלא כדי להבין אותו.

\subsection{מבוא להוכחה}

\L{%
\[
AB=\begin{pmatrix}a_{1,1} & \cdots & a_{1,n}\\
\vdots & \ddots & \vdots\\
a_{n,1} & \cdots & a_{n,n}
\end{pmatrix}\begin{pmatrix}b_{1,1} & \cdots & b_{1,n}\\
\vdots & \ddots & \vdots\\
b_{n,1} & \cdots & b_{n,n}
\end{pmatrix}=\begin{pmatrix}\sum_{k=1}^{n}a_{1,k}b_{k,1} & \cdots & \sum_{k=1}^{n}a_{1,k}b_{k,n}\\
\vdots & \ddots & \vdots\\
\sum_{k=1}^{n}a_{n,k}b_{k,1} & \cdots & \sum_{k=1}^{n}a_{n,k}b_{k,n}
\end{pmatrix}
\]
}%
\L{%
\[
=\begin{pmatrix}\sum_{k=1}^{n}a_{1,k}\frac{\left(-1\right)^{k+1}\det\left(A_{1,k}\right)}{\det(A)} & \cdots & \sum_{k=1}^{n}a_{1,k}\frac{\left(-1\right)^{k+n}\det\left(A_{n,k}\right)}{\det(A)}\\
\vdots & \ddots & \vdots\\
\sum_{k=1}^{n}a_{n,k}\frac{\left(-1\right)^{k+1}\det\left(A_{1,k}\right)}{\det(A)} & \cdots & \sum_{k=1}^{n}a_{n,k}\frac{\left(-1\right)^{k+n}\det\left(A_{n,k}\right)}{\det(A)}
\end{pmatrix}
\]
}%
\L{%
\[
=\frac{1}{\det(A)}\begin{pmatrix}\sum_{k=1}^{n}\left(-1\right)^{1+k}a_{1,k}\det\left(A_{1,k}\right) & \cdots & \sum_{k=1}^{n}\left(-1\right)^{n+k}a_{1,k}\det\left(A_{n,k}\right)\\
\vdots & \ddots & \vdots\\
\sum_{k=1}^{n}\left(-1\right)^{1+k}a_{n,k}\det\left(A_{1,k}\right) & \cdots & \sum_{k=1}^{n}\left(-1\right)^{n+k}a_{n,k}\det\left(A_{n,k}\right)
\end{pmatrix}
\]
}%
\L{%
\[
\underset{(1)}{=}\frac{1}{\det(A)}\begin{pmatrix}\det\begin{pmatrix}a_{1,1} & \cdots & a_{1,n}\\
a_{2,1} & \cdots & a_{2,n}\\
\vdots &  & \vdots\\
a_{n,1} & \cdots & a_{n,n}
\end{pmatrix} & \cdots & \det\begin{pmatrix}a_{1,1} & \cdots & a_{1,n}\\
\vdots &  & \vdots\\
a_{n-1,1} & \cdots & a_{n-1,n}\\
a_{1,1} & \cdots & a_{1,n}
\end{pmatrix}\\
\vdots & \ddots & \vdots\\
\det\begin{pmatrix}a_{n,1} & \cdots & a_{n,n}\\
a_{2,1} & \cdots & a_{2,n}\\
\vdots &  & \vdots\\
a_{n,1} & \cdots & a_{n,n}
\end{pmatrix} & \cdots & \det\begin{pmatrix}a_{1,1} & \cdots & a_{1,n}\\
\vdots &  & \vdots\\
a_{n-1,1} & \cdots & a_{n-1,n}\\
a_{n,1} & \cdots & a_{n,n}
\end{pmatrix}
\end{pmatrix}
\]
}%
\L{%
\[
=\frac{1}{\det(A)}\begin{pmatrix}\det(A) & \cdots & 0\\
\vdots & \ddots & \vdots\\
0 & \cdots & \det(A)
\end{pmatrix}=\begin{pmatrix}1 & 0 & \cdots & 0\\
0 & 1 & \cdots & 0\\
\vdots & \vdots & \ddots & \vdots\\
0 & 0 & \cdots & 1
\end{pmatrix}=I_{n}
\]
}%

\L{$(1)$} באלכסון )\L{$i=j$}(: מחליפים שורה \L{$j$} בעצמה, זו \L{$A$}
עצמה, והדטרמיננטה היא \L{$\det(A)$}.

מחוץ לאלכסון )\L{$i\ne j$}(: אותה שורה מופיעה פעמיים.

דטרמיננטה עם שתי שורות זהות היא \L{$0$}.\\

מאחר ש-\L{$AB=I_{n}$} ואלו מטריצות ריבועיות, \L{$BA=I_{n}$} )לינארית
1, הרצאה \L{$21$}(, מתקבל \L{$B=A^{-1}$}.

זו הנוסחה \L{$A^{-1}=\frac{1}{\det A}\text{adj}(A)$}.\\

רעיון ההוכחה - הקשר לפיתוח לפלס: פיתוח \L{$\det(A)$} לפי השורה ה-\L{$j$}
הוא \L{%
\[
\det(A)=\sum_{k=1}^{n}\left(-1\right)^{j+k}a_{j,k}\det\left(A_{j,k}\right)
\]
}%
כלומר כל איבר \L{$a_{j,k}$} בשורה \L{$j$} מוכפל ב-\L{$\left(-1\right)^{j+k}\det\left(A_{j,k}\right)$}.

האיבר במקום \L{$\left(j,j\right)$} של \L{$AB$} הוא שורה \L{$j$}
של \L{$A$} כפול עמודה \L{$j$} של \L{$B$}, ולפי הגדרת \L{$B$} \L{%
\[
\sum_{k=1}^{n}a_{j,k}b_{k,j}=\sum_{k=1}^{n}a_{j,k}\frac{\left(-1\right)^{j+k}\det\left(A_{j,k}\right)}{\det(A)}=\frac{1}{\det(A)}\underset{=\det(A)}{\underbrace{\sum_{k=1}^{n}\left(-1\right)^{j+k}a_{j,k}\det\left(A_{j,k}\right)}}=1
\]
}%
זה נותן את האלכסון. מחוץ לאלכסון מקבלים פיתוח של דטרמיננטה עם שתי
שורות זהות, ולכן \L{$0$}. זו ההוכחה שלהלן.
\begin{proof}
1. נסמן \L{%
\[
C=AB
\]
}%
\L{%
\[
C=\left(c_{i,j}\right)_{\substack{1\leq i\leq n\\
1\leq j\leq n
}
}
\]
}%

2. לפי הגדרת כפל מטריצות )לינארית 1, הרצאה \L{$19$}(, \L{%
\[
c_{i,j}=\sum_{\ell=1}^{n}a_{i,\ell}b_{\ell,j}
\]
}%
נציב את \L{$b_{\ell,j}$} )במקום האינדקס \L{$i$} של \L{$b_{i,j}$}
מופיע כעת אינדקס הסכימה \L{$\ell$}, ולכן מוחקים את העמודה ה-\L{$\ell$}
ואת השורה ה-\L{$j$}(: \L{%
\[
c_{i,j}=\sum_{\ell=1}^{n}a_{i,\ell}\frac{(-1)^{\ell+j}}{\det(A)}\det\left(\begin{array}{cc|c|cc}
a_{1,1} & \cdots & a_{1,\ell} & \cdots & a_{1,n}\\
\vdots &  & \vdots &  & \vdots\\
\hline a_{j,1} & \cdots & a_{j,\ell} & \cdots & a_{j,n}\\
\hline \vdots &  & \vdots &  & \vdots\\
a_{n,1} & \cdots & a_{n,\ell} & \cdots & a_{n,n}
\end{array}\right)\tag{\ensuremath{*}}
\]
}%
וזה נכון לכל \L{$1\leq i,j\leq n$}.

3. המקרה \L{$i=j$}: \L{%
\[
c_{i,i}=\frac{1}{\det(A)}\sum_{\ell=1}^{n}a_{i,\ell}(-1)^{\ell+i}\det\left(\begin{array}{cc|c|cc}
a_{1,1} & \cdots & a_{1,\ell} & \cdots & a_{1,n}\\
\vdots &  & \vdots &  & \vdots\\
\hline a_{i,1} & \cdots & a_{i,\ell} & \cdots & a_{i,n}\\
\hline \vdots &  & \vdots &  & \vdots\\
a_{n,1} & \cdots & a_{n,\ell} & \cdots & a_{n,n}
\end{array}\right)=\frac{\det(A)}{\det(A)}=1
\]
}%
כי הסכום הוא בדיוק פיתוח הדטרמיננטה של \L{$A$} לפי השורה ה-\L{$i$}.

לכן על האלכסון של \L{$C$} יש \L{$1$}-ים.

4. המקרה \L{$i\neq j$}. לצורך העניין נניח \L{$1\leq i<j\leq n$}
)המקרה \L{$i>j$} זהה(. נכפול את \L{$(*)$} ב-\L{$\det(A)$}: \L{%
\[
\det(A)\cdot c_{i,j}=\sum_{\ell=1}^{n}a_{i,\ell}(-1)^{\ell+j}\det\left(\begin{array}{cc|c|cc}
a_{1,1} & \cdots & a_{1,\ell} & \cdots & a_{1,n}\\
\vdots &  & \vdots &  & \vdots\\
a_{i,1} & \cdots & a_{i,\ell} & \cdots & a_{i,n}\\
\vdots &  & \vdots &  & \vdots\\
\hline a_{j,1} & \cdots & a_{j,\ell} & \cdots & a_{j,n}\\
\hline \vdots &  & \vdots &  & \vdots\\
a_{n,1} & \cdots & a_{n,\ell} & \cdots & a_{n,n}
\end{array}\right)
\]
}%
זהו בדיוק פיתוח לפי השורה ה-\L{$j$} של דטרמיננטה של מטריצה שבה במקום
השורה ה-\L{$j$} כתובה השורה ה-\L{$i$} של \L{$A$} )האיברים בשורה
זו הם \L{$a_{i,\ell}$}, והמינורים הם אותם מינורים, כי מחיקת השורה
ה-\L{$j$} מוחקת את השורה שהוחלפה(: \L{%
\[
\det(A)\cdot c_{i,j}=\det\begin{pmatrix}a_{1,1} & \cdots & a_{1,n}\\
\vdots &  & \vdots\\
a_{i,1} & \cdots & a_{i,n}\\
a_{i+1,1} & \cdots & a_{i+1,n}\\
\vdots &  & \vdots\\
a_{i,1} & \cdots & a_{i,n}\\
\vdots &  & \vdots\\
a_{n,1} & \cdots & a_{n,n}
\end{pmatrix}\;\leftarrow\text{row }j
\]
}%
במטריצה זו יש שתי שורות זהות )השורה ה-\L{$i$} והשורה ה-\L{$j$}(,
ולכן הדטרמיננטה שווה ל-\L{$0$} )לינארית 1, הרצאה \L{$26$}, טענה
11(.

מכאן \L{$\det(A)\cdot c_{i,j}=0$}, ומאחר ש-\L{$\det(A)\neq0_{\mathbb{F}}$},
\L{%
\[
c_{i,j}=0_{\mathbb{F}}
\]
}%

5. מ-3 ו-4: \L{$C=AB=I_{n}$}. 
\end{proof}
יבגני: זו למה טכנית, ונמשיך איתה ביום רביעי.
\end{document}
